\providecommand{\frR}{\mathfrak R}
\providecommand{\frS}{\mathfrak S}
\providecommand{\frG}{\mathfrak G}

\subsection*{\S 4. Systèmes hypercomplexes et leurs classes de représentations}

\paragraph{1. Théorème d'extension.}
Nous relions maintenant les théorèmes de représentation du \S 2 aux théorèmes sur les modules du \S 3. À la place des anneaux généraux \(\frR\), nous considérons des systèmes hypercomplexes avec élément unité -- \(S,T,\ldots\) -- sur un corps commutatif \(P\), donc
\[
        S=x_1P+\cdots+x_nP;
\]
à la place des anneaux de représentation arbitraires \(\frS\), seulement des corps \(A,B,\ldots\), et, sauf remarque contraire, de tels corps avec centre \(P\). Dans ce cas il existe toujours l'anneau-produit dans lequel \(S\) et \(A\) deviennent permutables élément par élément, avec intersection \(P\); il s'agit du produit direct \(S\times A\), qui devient en même temps le module d'extension \(S_A\) de \(S\) au sens du \S 3,
\[
        S_A=x_1A+\cdots+x_nA,
\]
et qui est donc aussi désigné comme anneau d'extension.

Dans cette spécialisation, le théorème sur les modules invariants (\S 3, 3.) devient le

\paragraph{Théorème d'extension.}
Tout idéal bilatéral de \(S_A\) est idéal d'extension d'un idéal de \(S\), à savoir extension de son intersection avec \(S\).

En effet, si le groupe \(\frG\) pris pour base est celui des automorphismes intérieurs de \(A\), alors \(P\), centre de \(A\), devient corps complet des invariants, tandis que la permutabilité élément par élément de \(A\) avec \(S\) dit que \(\frG\), d'après \S 3, 3., est étendu à \(S_A\) de telle sorte que \(S\) devienne domaine complet des invariants. Tout idéal bilatéral \(\mathfrak a\) de \(S_A\) est un sous-groupe admissible -- à cause de \(\alpha^{-1}\mathfrak a\alpha\subseteq\mathfrak a\) -- donc, d'après \S 3, 3., extension de son module-intersection \(\mathfrak a\cap S\), qui devient un idéal de \(S\).

Addendum: Si \(S\) est commutatif, alors \(S\) devient le centre de \(S_A\). Car \(S\) est contenu dans le centre et est domaine complet des invariants par rapport aux automorphismes intérieurs engendrés par \(A\). Plus généralement, le centre de \(S\) devient aussi le centre de \(S_A\).

Notation: Par \(A_r,B_n,\ldots\) on entendra toujours les anneaux matriciels du degré indiqué sur \(A,B,\ldots\), donc
\[
        A_r=\sum A c_{ik}=\sum c_{ik}A.
\]

Dans ce qui suit on utilisera essentiellement une conséquence, le théorème d'extension pour les systèmes simples: si \(S\) est un système simple\footnote{Système simple signifie, comme d'habitude, ``simple bilatéral''.}, alors \(S_A\) devient lui aussi simple:
\[
        S_A=\sum B c_{ik}=\sum c_{ik}B=B_t
\]
avec corps \(B\) associé, dont le centre contient \(P\). Ici \(A\), et donc aussi \(B\), peut être de rang infini sur \(P\); seul \(S\) est supposé hypercomplexe. Si \(S\) et \(A\) possèdent le centre commun \(P\), alors \(P\) devient aussi centre de \(S_A\).

Plus généralement encore: si \(S\) est un système hypercomplexe simple, alors le produit direct \(S\times A_r\) devient lui aussi simple. Car \(S\times A_r\) coïncide avec \(S_A\times P_r\), et devient donc égal à \(B_{tr}\), lorsque \(S_A=B_t\).

\paragraph{2. Classes de représentations.}
Par représentation -- réciproque ou directe -- d'un système hypercomplexe, on entendra, comme d'habitude, une représentation homomorphe relativement aux opérateurs du domaine de coefficients \(P\) (\S 2, fin) et distincte de la représentation nulle. La liaison du théorème de passage et de sa conséquence (\S 2, 3.) avec le théorème d'extension de 1. donne le

\paragraph{Théorème sur les classes de représentations.}
Si \(S\) est hypercomplexe avec élément unité, avec \(P\) comme domaine de coefficients, et si \(A\) est un corps de centre \(P\), alors \(S\) possède autant de classes irréductibles différentes de représentations réciproques de \(S\) dans \(A\) qu'il y a de classes d'idéaux à droite simples dans l'anneau quotient de \(S_A\) par son radical. -- Si \(S\) est un système simple, alors il possède exactement une classe irréductible de représentations réciproques et une classe irréductible de représentations directes dans \(A\).\footnote{Voir Darstellungstheorie, note 20), pour le cas où \(A\) est le corps gauche des endomorphismes associé de \(S\).}

Car, d'après la conséquence du théorème de passage (\S 2, 3.), les classes simples de représentations réciproques et les classes simples de modules sur \(S_A\) se correspondent biunivoquement. Puisque \(S_A\) est de rang fini sur le corps \(A\), il satisfait aux conditions maximale et minimale pour les idéaux unilatères; ainsi la première partie de l'énoncé découle directement de Darstellungstheorie, \S 19. Si \(S\) est simple, alors \(S_A\) l'est aussi d'après le théorème d'extension. À cause des conditions maximale et minimale, il est en même temps complètement réductible, donc ne possède qu'une seule classe d'idéaux unilatères simples, c'est-à-dire que \(S\) ne possède qu'une seule classe irréductible de représentations réciproques dans \(A\). Avec \(S\), le système qui lui est anti-isomorphe est lui aussi simple et ne possède qu'une seule classe irréductible de représentations réciproques; donc \(S\) ne possède aussi qu'une seule classe simple de représentations directes.

\paragraph{3. Relation de rang.}
Le fait que, si \(S\) est simple,
\[
        S_A=\sum Bc_{ik}
\]
est de rang fini aussi bien sur \(A\) que sur \(B\), donne la relation de rang:
\[
        n=rt\quad\text{avec}\quad n=(S:P),\quad t^2=(S_A:B),
        \quad r=\text{degré de la représentation réciproque irréductible}.      \tag{1}
\]

En effet, \(S_A\) est lui-même module de représentation réciproque pour \(S\) dans \(A\) (la représentation se trouve déjà dans \(P\)). Comme anneau matriciel sur \(B\), \(S_A\) se décompose en \(t\) idéaux à droite simples, donc en \(S_A\)-modules, dont le rang sur \(A\) est égal au degré \(r\) de la représentation réciproque irréductible. Puisque \(n\) est en même temps le rang de \(S_A\) sur \(A\), la relation de rang en résulte.

\paragraph{4. Renforcement de la relation de rang, lorsque \(A\) est de rang fini sur \(P\).}
Dans ce cas les trois relations valent:
\[
        (S:P)=n=rt,                                                     \tag{1}
\]
\[
        (B:P)\,t=(A:P)\,r,                                               \tag{2}
\]
\[
        (S:P)(B:P)=(A:P)\,r^2.                                           \tag{3}
\]
Ici (2) signifie le rang sur \(P\) d'un idéal à droite simple de \(S_A\), exprimé d'après 3. par le rang sur \(B\), respectivement sur \(A\), tandis que (3) signifie le rang d'un idéal à droite simple dans un anneau matriciel \(B_n\) et résulte de (2) par multiplication par \(r\), en vertu de (1).

\clearpage
